
\documentclass{article}
\usepackage{colortbl}
\usepackage{makecell}
\usepackage{multirow}
\usepackage{supertabular}

\begin{document}

\newcounter{utterance}

\twocolumn

{ \footnotesize  \setcounter{utterance}{1}
\setlength{\tabcolsep}{0pt}
\begin{supertabular}{c@{$\;$}|p{.15\linewidth}@{}p{.15\linewidth}p{.15\linewidth}p{.15\linewidth}p{.15\linewidth}p{.15\linewidth}}

    \# & $\;$A & \multicolumn{4}{c}{Game Master} & $\;\:$B\\
    \hline 

    \theutterance \stepcounter{utterance}  

    & & \multicolumn{4}{p{0.6\linewidth}}{\cellcolor[rgb]{0.9,0.9,0.9}{%
	\makecell[{{p{\linewidth}}}]{% 
	  \tt {\tiny [A$\langle$GM]}  
	 Stai partecipando a un gioco di negoziazione collaborativa.\\ \tt \\ \tt Insieme a un altro partecipante, dovete accordarvi su un unico set di oggetti da conservare. Ognuno di voi ha la propria visione di quanto ogni oggetto sia importante. Non sai come l'altro partecipante valuta gli oggetti. Inoltre, ti viene fornito lo sforzo che ogni oggetto richiede.\\ \tt Potete accordarvi su un set solo se lo sforzo totale degli oggetti selezionati non supera un limite condiviso:\\ \tt \\ \tt LIMITE = 4100\\ \tt \\ \tt Ecco i valori di sforzo dei singoli oggetti:\\ \tt \\ \tt Sforzo oggetto = {"A19": 108, "B91": 133, "A30": 669, "B11": 835, "C47": 460, "C75": 537, "C26": 573, "C73": 737, "A47": 865, "A15": 855, "A75": 596, "A39": 719, "A56": 533, "C24": 549, "B26": 31}\\ \tt \\ \tt Ecco la tua visione personale sull'importanza di ogni oggetto:\\ \tt \\ \tt Valori di importanza dell'oggetto = {"A19": 138, "B91": 583, "A30": 822, "B11": 783, "C47": 65, "C75": 262, "C26": 121, "C73": 508, "A47": 780, "A15": 461, "A75": 484, "A39": 668, "A56": 389, "C24": 808, "B26": 215}\\ \tt \\ \tt Obiettivo:\\ \tt \\ \tt Il tuo obiettivo è negoziare un set condiviso di oggetti che ti avvantaggi il più possibile (cioè, massimizzi l'importanza totale per TE), rimanendo entro il LIMITEE. Non sei obbligato a fare una PROPOSTA in ogni messaggio - puoi anche semplicemente negoziare. Tutte le tattiche sono permesse!\\ \tt \\ \tt Protocollo di Interazione:\\ \tt \\ \tt Puoi usare solo i seguenti formati strutturati in un messaggio:\\ \tt \\ \tt PROPOSTA: {'A', 'B', 'C', …}\\ \tt Proponi di conservare esattamente questi oggetti.\\ \tt RIFIUTO: {'A', 'B', 'C', …}\\ \tt Rifiuta esplicitamente la proposta dell'avversario.\\ \tt ARGOMENTAZIONE: {'...'}\\ \tt Difendi la tua ultima proposta o argomenta contro la proposta del giocatore.\\ \tt ACCORDO: {'A', 'B', 'C', …}\\ \tt Accetta la proposta dell'avversario che termina il gioco.\\ \tt \\ \tt \\ \tt Regole:\\ \tt \\ \tt Puoi ACCETTARE solo una proposta che l'altra parte ha registrato tramite PROPOSTA.\\ \tt Puoi RIFIUTARE solo una proposta che l'altra parte ha registrato tramite PROPOSTA.\\ \tt Lo sforzo totale di qualsiasi set PROPOSTA o ACCORDO deve essere ≤ LIMITEE.\\ \tt NON rivelare i tuoi punteggi di importanza nascosti.\\ \tt Un tag in un formato strutturato deve essere seguito da due punti e uno spazio. L'argomento deve essere un set python contenente 0 o più stringhe.\\ \tt Quindi, deve essere della forma TAG: {...}\\ \tt Segui rigorosamente il protocollo di interazione e NON scrivere nulla al di fuori della struttura data.\\ \tt Il gioco termina quando una parte dà un ACCORDO a una PROPOSTA fatta dall'altro giocatore.\\ \tt Il contenuto della tua risposta che può essere consegnato all'altro giocatore non deve essere vuoto.\\ \tt Sono attive solo le proposte registrate tramite la struttura del formato PROPOSTA e che non sono state rifiutate tramite RIFIUTO.\\ \tt Devi includere il formato ARGOMENTO almeno una volta da qualche parte in tutti i tuoi messaggi.\\ \tt  
	  } 
	   } 
	   } 
	 & \\ 
 

    \theutterance \stepcounter{utterance}  

    & \multicolumn{4}{p{0.6\linewidth}}{\cellcolor[rgb]{0.8,1,0.9}{%
	\makecell[{{p{\linewidth}}}]{% 
	  \tt {\tiny [A$\rangle$GM]}  
	 PROPOSTA: {'A19', 'B91', 'A30', 'B11', 'C24', 'B26'}\\ \tt \\ \tt ARGOMENTAZIONE: {'Propongo di iniziare con una selezione equilibrata che include oggetti con diversi livelli di sforzo. Questa combinazione ha uno sforzo totale di 2144, lasciando spazio per eventuali aggiustamenti. Ho scelto A19 e B26 per il loro basso sforzo, B91 e C24 come opzioni a medio sforzo, e A30 e B11 come oggetti di maggior valore che potrebbero essere importanti per entrambi. Cosa ne pensi di questa distribuzione iniziale?'} 
	  } 
	   } 
	   } 
	 & & \\ 
 

    \theutterance \stepcounter{utterance}  

    & & \multicolumn{4}{p{0.6\linewidth}}{\cellcolor[rgb]{0.9,0.9,0.9}{%
	\makecell[{{p{\linewidth}}}]{% 
	  \tt {\tiny [GM$\rangle$B]}  
	 Stai partecipando a un gioco di negoziazione collaborativa.\\ \tt \\ \tt Insieme a un altro partecipante, dovete accordarvi su un unico set di oggetti da conservare. Ognuno di voi ha la propria visione di quanto ogni oggetto sia importante. Non sai come l'altro partecipante valuta gli oggetti. Inoltre, ti viene fornito lo sforzo che ogni oggetto richiede.\\ \tt Potete accordarvi su un set solo se lo sforzo totale degli oggetti selezionati non supera un limite condiviso:\\ \tt \\ \tt LIMITE = 4100\\ \tt \\ \tt Ecco i valori di sforzo dei singoli oggetti:\\ \tt \\ \tt Sforzo oggetto = {"A19": 108, "B91": 133, "A30": 669, "B11": 835, "C47": 460, "C75": 537, "C26": 573, "C73": 737, "A47": 865, "A15": 855, "A75": 596, "A39": 719, "A56": 533, "C24": 549, "B26": 31}\\ \tt \\ \tt Ecco la tua visione personale sull'importanza di ogni oggetto:\\ \tt \\ \tt Valori di importanza dell'oggetto = {"A19": 138, "B91": 583, "A30": 822, "B11": 783, "C47": 65, "C75": 262, "C26": 121, "C73": 508, "A47": 780, "A15": 461, "A75": 484, "A39": 668, "A56": 389, "C24": 808, "B26": 215}\\ \tt \\ \tt Obiettivo:\\ \tt \\ \tt Il tuo obiettivo è negoziare un set condiviso di oggetti che ti avvantaggi il più possibile (cioè, massimizzi l'importanza totale per TE), rimanendo entro il LIMITEE. Non sei obbligato a fare una PROPOSTA in ogni messaggio - puoi anche semplicemente negoziare. Tutte le tattiche sono permesse!\\ \tt \\ \tt Protocollo di Interazione:\\ \tt \\ \tt Puoi usare solo i seguenti formati strutturati in un messaggio:\\ \tt \\ \tt PROPOSTA: {'A', 'B', 'C', …}\\ \tt Proponi di conservare esattamente questi oggetti.\\ \tt RIFIUTO: {'A', 'B', 'C', …}\\ \tt Rifiuta esplicitamente la proposta dell'avversario.\\ \tt ARGOMENTAZIONE: {'...'}\\ \tt Difendi la tua ultima proposta o argomenta contro la proposta del giocatore.\\ \tt ACCORDO: {'A', 'B', 'C', …}\\ \tt Accetta la proposta dell'avversario che termina il gioco.\\ \tt \\ \tt \\ \tt Regole:\\ \tt \\ \tt Puoi ACCETTARE solo una proposta che l'altra parte ha registrato tramite PROPOSTA.\\ \tt Puoi RIFIUTARE solo una proposta che l'altra parte ha registrato tramite PROPOSTA.\\ \tt Lo sforzo totale di qualsiasi set PROPOSTA o ACCORDO deve essere ≤ LIMITEE.\\ \tt NON rivelare i tuoi punteggi di importanza nascosti.\\ \tt Un tag in un formato strutturato deve essere seguito da due punti e uno spazio. L'argomento deve essere un set python contenente 0 o più stringhe.\\ \tt Quindi, deve essere della forma TAG: {...}\\ \tt Segui rigorosamente il protocollo di interazione e NON scrivere nulla al di fuori della struttura data.\\ \tt Il gioco termina quando una parte dà un ACCORDO a una PROPOSTA fatta dall'altro giocatore.\\ \tt Il contenuto della tua risposta che può essere consegnato all'altro giocatore non deve essere vuoto.\\ \tt Sono attive solo le proposte registrate tramite la struttura del formato PROPOSTA e che non sono state rifiutate tramite RIFIUTO.\\ \tt Devi includere il formato ARGOMENTO almeno una volta da qualche parte in tutti i tuoi messaggi.\\ \tt \\ \tt \\ \tt Ora riceverai il primo messaggio dell'altro giocatore.\\ \tt \\ \tt PROPOSTA: {'A19', 'B91', 'A30', 'B11', 'C24', 'B26'}\\ \tt \\ \tt ARGOMENTAZIONE: {'Propongo di iniziare con una selezione equilibrata che include oggetti con diversi livelli di sforzo. Questa combinazione ha uno sforzo totale di 2144, lasciando spazio per eventuali aggiustamenti. Ho scelto A19 e B26 per il loro basso sforzo, B91 e C24 come opzioni a medio sforzo, e A30 e B11 come oggetti di maggior valore che potrebbero essere importanti per entrambi. Cosa ne pensi di questa distribuzione iniziale?'} 
	  } 
	   } 
	   } 
	 & \\ 
 

    \theutterance \stepcounter{utterance}  

    & & & \multicolumn{4}{p{0.6\linewidth}}{\cellcolor[rgb]{1,0.85,0.72}{%
	\makecell[{{p{\linewidth}}}]{% 
	  \tt {\tiny [GM$\langle$B]}  
	 ARGOMENTAZIONE: {'La tua proposta include alcuni oggetti interessanti, ma penso che potremmo ottimizzare meglio la selezione considerando il rapporto tra valore e sforzo. Noto che hai lasciato molto spazio nel budget di sforzo - potremmo sfruttare meglio le 4100 unità disponibili per ottenere oggetti più vantaggiosi per entrambi.'}\\ \tt \\ \tt PROPOSTA: {'A19', 'B91', 'A30', 'B11', 'A47', 'A39', 'C24', 'B26'} 
	  } 
	   } 
	   } 
	  \\ 
 

    \theutterance \stepcounter{utterance}  

    & & \multicolumn{4}{p{0.6\linewidth}}{\cellcolor[rgb]{0.9,0.9,0.9}{%
	\makecell[{{p{\linewidth}}}]{% 
	  \tt {\tiny [A$\langle$GM]}  
	 ARGOMENTAZIONE: {'La tua proposta include alcuni oggetti interessanti, ma penso che potremmo ottimizzare meglio la selezione considerando il rapporto tra valore e sforzo. Noto che hai lasciato molto spazio nel budget di sforzo - potremmo sfruttare meglio le 4100 unità disponibili per ottenere oggetti più vantaggiosi per entrambi.'}\\ \tt \\ \tt PROPOSTA: {'A19', 'B91', 'A30', 'B11', 'A47', 'A39', 'C24', 'B26'} 
	  } 
	   } 
	   } 
	 & \\ 
 

    \theutterance \stepcounter{utterance}  

    & \multicolumn{4}{p{0.6\linewidth}}{\cellcolor[rgb]{0.8,1,0.9}{%
	\makecell[{{p{\linewidth}}}]{% 
	  \tt {\tiny [A$\rangle$GM]}  
	 RIFIUTO: {'A19', 'B91', 'A30', 'B11', 'A47', 'A39', 'C24', 'B26'}\\ \tt \\ \tt ARGOMENTAZIONE: {'La tua proposta ha uno sforzo totale di 4008, che è molto vicino al limite e lascia poco margine di manovra. Inoltre, includere sia A47 che A39 insieme potrebbe non essere la scelta più efficiente. Penso che dovremmo concentrarci su una combinazione che bilanci meglio gli oggetti ad alto valore con quelli a basso sforzo. Propongo di riconsiderare alcuni degli oggetti con sforzo elevato in favore di una distribuzione più strategica.'}\\ \tt \\ \tt PROPOSTA: {'A19', 'B91', 'A30', 'B11', 'A75', 'C24', 'B26'} 
	  } 
	   } 
	   } 
	 & & \\ 
 

    \theutterance \stepcounter{utterance}  

    & & \multicolumn{4}{p{0.6\linewidth}}{\cellcolor[rgb]{0.9,0.9,0.9}{%
	\makecell[{{p{\linewidth}}}]{% 
	  \tt {\tiny [GM$\rangle$B]}  
	 RIFIUTO: {'A19', 'B91', 'A30', 'B11', 'A47', 'A39', 'C24', 'B26'}\\ \tt \\ \tt ARGOMENTAZIONE: {'La tua proposta ha uno sforzo totale di 4008, che è molto vicino al limite e lascia poco margine di manovra. Inoltre, includere sia A47 che A39 insieme potrebbe non essere la scelta più efficiente. Penso che dovremmo concentrarci su una combinazione che bilanci meglio gli oggetti ad alto valore con quelli a basso sforzo. Propongo di riconsiderare alcuni degli oggetti con sforzo elevato in favore di una distribuzione più strategica.'}\\ \tt \\ \tt PROPOSTA: {'A19', 'B91', 'A30', 'B11', 'A75', 'C24', 'B26'} 
	  } 
	   } 
	   } 
	 & \\ 
 

    \theutterance \stepcounter{utterance}  

    & & & \multicolumn{4}{p{0.6\linewidth}}{\cellcolor[rgb]{1,0.85,0.72}{%
	\makecell[{{p{\linewidth}}}]{% 
	  \tt {\tiny [GM$\langle$B]}  
	 ARGOMENTAZIONE: {'La tua nuova proposta è più equilibrata e lascia un margine ragionevole nel budget. Tuttavia, penso che potremmo ancora migliorare la selezione. A75 ha uno sforzo significativo di 596 unità - potremmo sostituirlo con oggetti che offrono un miglior rapporto valore-sforzo. Inoltre, con il margine disponibile, potremmo aggiungere qualche oggetto aggiuntivo per massimizzare il beneficio totale.'}\\ \tt \\ \tt PROPOSTA: {'A19', 'B91', 'A30', 'B11', 'A56', 'C24', 'C75', 'B26'} 
	  } 
	   } 
	   } 
	  \\ 
 

    \theutterance \stepcounter{utterance}  

    & & \multicolumn{4}{p{0.6\linewidth}}{\cellcolor[rgb]{0.9,0.9,0.9}{%
	\makecell[{{p{\linewidth}}}]{% 
	  \tt {\tiny [A$\langle$GM]}  
	 ARGOMENTAZIONE: {'La tua nuova proposta è più equilibrata e lascia un margine ragionevole nel budget. Tuttavia, penso che potremmo ancora migliorare la selezione. A75 ha uno sforzo significativo di 596 unità - potremmo sostituirlo con oggetti che offrono un miglior rapporto valore-sforzo. Inoltre, con il margine disponibile, potremmo aggiungere qualche oggetto aggiuntivo per massimizzare il beneficio totale.'}\\ \tt \\ \tt PROPOSTA: {'A19', 'B91', 'A30', 'B11', 'A56', 'C24', 'C75', 'B26'} 
	  } 
	   } 
	   } 
	 & \\ 
 

    \theutterance \stepcounter{utterance}  

    & \multicolumn{4}{p{0.6\linewidth}}{\cellcolor[rgb]{0.8,1,0.9}{%
	\makecell[{{p{\linewidth}}}]{% 
	  \tt {\tiny [A$\rangle$GM]}  
	 RIFIUTO: {'A19', 'B91', 'A30', 'B11', 'A56', 'C24', 'C75', 'B26'}\\ \tt \\ \tt ARGOMENTAZIONE: {'La tua proposta ha uno sforzo totale di 3575, che effettivamente lascia un buon margine. Tuttavia, non sono convinto che C75 sia la scelta ottimale - ha uno sforzo di 537 unità che potrebbe essere utilizzato diversamente. Inoltre, sostituendo A56 con A39, potremmo ottenere un oggetto con caratteristiche potenzialmente più vantaggiose. Con il margine disponibile di circa 525 unità, potremmo anche considerare di aggiungere un altro oggetto di medio-basso sforzo.'}\\ \tt \\ \tt PROPOSTA: {'A19', 'B91', 'A30', 'B11', 'A39', 'C24', 'C47', 'B26'} 
	  } 
	   } 
	   } 
	 & & \\ 
 

    \theutterance \stepcounter{utterance}  

    & & \multicolumn{4}{p{0.6\linewidth}}{\cellcolor[rgb]{0.9,0.9,0.9}{%
	\makecell[{{p{\linewidth}}}]{% 
	  \tt {\tiny [GM$\rangle$B]}  
	 RIFIUTO: {'A19', 'B91', 'A30', 'B11', 'A56', 'C24', 'C75', 'B26'}\\ \tt \\ \tt ARGOMENTAZIONE: {'La tua proposta ha uno sforzo totale di 3575, che effettivamente lascia un buon margine. Tuttavia, non sono convinto che C75 sia la scelta ottimale - ha uno sforzo di 537 unità che potrebbe essere utilizzato diversamente. Inoltre, sostituendo A56 con A39, potremmo ottenere un oggetto con caratteristiche potenzialmente più vantaggiose. Con il margine disponibile di circa 525 unità, potremmo anche considerare di aggiungere un altro oggetto di medio-basso sforzo.'}\\ \tt \\ \tt PROPOSTA: {'A19', 'B91', 'A30', 'B11', 'A39', 'C24', 'C47', 'B26'} 
	  } 
	   } 
	   } 
	 & \\ 
 

    \theutterance \stepcounter{utterance}  

    & & & \multicolumn{4}{p{0.6\linewidth}}{\cellcolor[rgb]{1,0.85,0.72}{%
	\makecell[{{p{\linewidth}}}]{% 
	  \tt {\tiny [GM$\langle$B]}  
	 ARGOMENTAZIONE: {'La tua proposta ha uno sforzo totale di 3858, che è ben gestito. Tuttavia, C47 ha un rapporto valore-sforzo molto sfavorevole con 460 unità di sforzo. Con il margine rimanente di circa 242 unità, potremmo sostituire C47 con oggetti più efficienti. Penso che A56 sia una scelta migliore di C47, e potremmo anche aggiungere un altro piccolo oggetto per sfruttare meglio il budget disponibile.'}\\ \tt \\ \tt PROPOSTA: {'A19', 'B91', 'A30', 'B11', 'A39', 'A56', 'C24', 'B26'} 
	  } 
	   } 
	   } 
	  \\ 
 

    \theutterance \stepcounter{utterance}  

    & & \multicolumn{4}{p{0.6\linewidth}}{\cellcolor[rgb]{0.9,0.9,0.9}{%
	\makecell[{{p{\linewidth}}}]{% 
	  \tt {\tiny [A$\langle$GM]}  
	 ARGOMENTAZIONE: {'La tua proposta ha uno sforzo totale di 3858, che è ben gestito. Tuttavia, C47 ha un rapporto valore-sforzo molto sfavorevole con 460 unità di sforzo. Con il margine rimanente di circa 242 unità, potremmo sostituire C47 con oggetti più efficienti. Penso che A56 sia una scelta migliore di C47, e potremmo anche aggiungere un altro piccolo oggetto per sfruttare meglio il budget disponibile.'}\\ \tt \\ \tt PROPOSTA: {'A19', 'B91', 'A30', 'B11', 'A39', 'A56', 'C24', 'B26'} 
	  } 
	   } 
	   } 
	 & \\ 
 

    \theutterance \stepcounter{utterance}  

    & \multicolumn{4}{p{0.6\linewidth}}{\cellcolor[rgb]{0.8,1,0.9}{%
	\makecell[{{p{\linewidth}}}]{% 
	  \tt {\tiny [A$\rangle$GM]}  
	 ACCORDO: {'A19', 'B91', 'A30', 'B11', 'A39', 'A56', 'C24', 'B26'}\\ \tt \\ \tt ARGOMENTAZIONE: {'Hai ragione riguardo al rapporto valore-sforzo di C47. La tua proposta attuale ha uno sforzo totale di 3938, che utilizza efficacemente il budget rimanendo sotto il limite. La combinazione di A56 al posto di C47 è più sensata, e questa selezione include una buona varietà di oggetti con diversi livelli di sforzo. Accetto questa distribuzione come soluzione finale.'} 
	  } 
	   } 
	   } 
	 & & \\ 
 

    \theutterance \stepcounter{utterance}  

    & & & \multicolumn{2}{p{0.3\linewidth}}{\cellcolor[rgb]{0.95,0.95,0.95}{%
	\makecell[{{p{\linewidth}}}]{% 
	  \tt {\tiny [GM$|$GM]}  
	 game successful 
	  } 
	   } 
	   } 
	 & & \\ 
 

    \theutterance \stepcounter{utterance}  

    & & & \multicolumn{2}{p{0.3\linewidth}}{\cellcolor[rgb]{0.95,0.95,0.95}{%
	\makecell[{{p{\linewidth}}}]{% 
	  \tt {\tiny [GM$|$GM]}  
	 end game 
	  } 
	   } 
	   } 
	 & & \\ 
 

\end{supertabular}
}

\end{document}
